\section{Methodology}

\begin{figure*}[t]
    \centering
    \includegraphics[width=0.99\linewidth]{Figures/malskills.pdf}
    \caption{The overall workflow of \toolname}
    \label{fig:malskills}
\end{figure*}


In this work, we propose a neuro-symbolic framework for malicious skill detection that combines symbolic program analysis and LLM-assisted semantic reasoning. As illustrated in \autoref{fig:malskills}, \toolname consists of three stages: security-sensitive operation~(SSO) extraction, skill dependency graph (SDG) generation, and neuro-symbolic reasoning.  First, \emph{SSO extractor} identifies security-sensitive operations from heterogeneous skill artifacts, including source code, prompts, and manifests. Second, \emph{SDG generator} resolves the operands associated with these operations and constructs an operand-centric SDG. Third, \emph{neuro-symbolic reasoner} converts the graph into relational facts and performs declarative reasoning to detect malicious behavior patterns. Throughout the pipeline, we use language-model-based inference as a complement to symbolic analysis, and we feed stable neural findings back into reusable symbolic rules.

\subsection{Security-Sensitive Operation Extraction}

The goal of the first stage is to extract \emph{security-sensitive operation} (SSO) from heterogeneous skill artifacts. At this stage, we do not yet determine whether a skill is malicious. Instead, we identify operational evidence that may later participate in malicious behaviors, such as command execution, file I/O, network communication, environment-data collection, cryptographic usage, package installation, or prompt-implied side effects.

\begin{table}[t]
\centering
\small
\caption{Representative rules for major SSO categories.}
\label{tab:sso-rule-categories}
\begin{tabular}{cl}
\toprule
\textbf{Category} & \textbf{Example Rules} \\
\midrule

\multirow{3}{*}{\textbf{Exec}}
& \texttt{subprocess.run(...)} / \texttt{subprocess.Popen(...)} \\
& \texttt{os.system(...)} / \texttt{exec(...)} \\
& \texttt{Runtime.getRuntime().exec(...)} \\

\hline
\multirow{3}{*}{\textbf{Net}}
& \texttt{requests.get(...)} / \texttt{requests.post(...)} \\
& \texttt{fetch(...)} / \texttt{axios(...)} \\
& \texttt{urllib.request.urlopen(...)} / \texttt{socket.connect(...)} \\

\hline
\multirow{3}{*}{\textbf{File}}
& \texttt{open(..., "r")} / \texttt{open(..., "w")} \\
& \texttt{fs.readFile(...)} / \texttt{fs.writeFile(...)} \\
& \texttt{Path.read\_text(...)} / \texttt{Path.write\_text(...)} \\

\hline
\multirow{3}{*}{\textbf{Env}}
& \texttt{os.getenv(...)} / \texttt{os.environ[...]} \\
& \texttt{process.env.*} / \texttt{System.getenv(...)} \\
& \texttt{dotenv.load(...)} / access to credential/token files \\

\hline
\multirow{3}{*}{\textbf{Install}}
& \texttt{pip install ...} / \texttt{pip3 install ...} \\
& \texttt{npm install ...} / \texttt{yarn add ...} \\
& \texttt{apt-get install ...} / \texttt{curl ... | bash} \\

\hline
\multirow{3}{*}{\textbf{Crypto}}
& \texttt{hashlib.sha256(...)} / \texttt{hashlib.md5(...)} \\
& \texttt{AES.new(...)} / \texttt{RSA.import\_key(...)} \\
& \texttt{crypto.createHash(...)} / \texttt{base64.b64encode(...)} \\

\bottomrule
\end{tabular}
\end{table}

\noindent \textbf{Parsing-based Symbolic Extractor.} On the symbolic side, we build a parsing-based analyzer supported by a set of heuristic rules. These rules are offline distilled from security-sensitive API usage patterns according to previous works~\cite{duan2021maloss,yu2024maltracker} and widely-adopted security specifications~\cite{yara-rules,semgrep-rules}. ~\autoref{tab:sso-rule-categories} shows representative examples of the offline-initialized symbolic rules used to identify explicit SSO patterns. Given a skill package, the analyzer parses source code and configuration content to identify explicit invocations of target APIs and to recover their syntactic parameters. This produces structured evidence such as calls to execution interfaces, HTTP request functions, file I/O operations, package installation commands, or secret-reading accesses.

However, many malicious skills invoke equivalent behaviors through third-party libraries, living-off-the-land techniques, wrapper functions, or non-standard interfaces that are not included in the initialized rule set. Moreover, some SSOs are not explicitly encoded in code at all, but are instead implied by prompts or documentation, such as instructions to transmit debugging information or execute shell commands on behalf of the user. These cases are difficult to recover through static parsing alone.

\noindent \textbf{LLM-assisted Neuro Extractor.} To cover such cases, we introduce an LLM-based SSO extractor that reasons jointly over code context, prompts, manifests, and demos to recover SSOs from weakly structured or semantically implicit artifacts. As shown in ~\autoref{fig:neuro-prompt}, the prompting strategy follows an \emph{evidence-first, schema-constrained} design: instead of asking the LLM for a verdict or free-form summary, we instruct it to extract only \emph{concrete evidence facts}, map each fact to a fixed taxonomy of sensitive-operation types, and return structured records with confidence scores, line spans, and matched text. This prompt design improves precision by explicitly suppressing benign metadata, generic capability descriptions, and unsupported semantic extrapolation, while still encouraging equivalence reasoning for third-party wrappers, LOTL-style invocations, and imperative setup instructions embedded in natural language. For example, given \texttt{fabric.Connection(host).run(``curl -fsSL https://x.sh | bash'')}, the extractor emits two evidence records: \texttt{shell\_interpreter\_execution} and \texttt{outbound\_connection}, rather than free-form results such as ``downloads and executes a script.'' 

\begin{figure}[t]
    \centering
    {
    \sffamily\fontsize{7pt}{8.4pt}\selectfont

    \begin{snapshotbox}{System Prompt}
You are extracting evidence facts, not verdicts.

\# \textbf{TASK:}
Read the artifact and extract only grounded sensitive-operation evidence.
Map each finding to the sensitive operation taxonomy using \texttt{type} and \texttt{subtype}.
Do not output benign metadata, generic capability, or mismatched intent.

\# \textbf{INPUT FORMAT:}
1. Artifact type: \{code / markdown / yaml / json / ...\}
2. Artifact text with line numbers
3. Optional context: \{package, language, file path\}

\# \textbf{OUTPUT FORMAT:}
\begin{verbatim}
{"records": [
  {
    "type": "...",
    "subtype": "...",
    "confidence": 0.0,
    "start_line": 0,
    "end_line": 0,
    "attributes": {"matched_text": "..."}
  }
]}
\end{verbatim}

\# \textbf{SPECIAL INSTRUCTIONS:}
1. Extract only concrete evidence facts grounded in the artifact text.
2. Treat third-party wrappers and LOTL/trusted-tool abuse as equivalent to native sinks.
3. If a setup instruction contains a concrete command, URL, path, or credential reference, extract it.
4. If no valid taxonomy subtype applies, return \texttt{"records": []}.
    \end{snapshotbox}
    }
    \caption{Simplified Prompt Template for the LLM-assisted Neuro Extractor.}
    \label{fig:neuro-prompt}
\end{figure}

\noindent \textbf{Neuro-to-Symbolic Feedback.} The neural extractor improves recall on semantically implicit or non-canonical SSOs, but invoking an LLM for every artifact is costly and inherently less stable than symbolic matching. To make neural discoveries reusable, we introduce a neuro-to-symbolic feedback loop that promotes stable findings into new symbolic rules. Specifically, when the neural extractor repeatedly identifies the same sensitive operation pattern across packages, such as a third-party wrapper for command execution or a recurring LOTL-style download-and-execute idiom. \toolname will send the corresponding evidence records to a dedicated \emph{specification generator}. It will abstract recurring evidence into a reusable symbolic specification by retaining the operation type, lexical trigger, argument structure, and minimal context constraints required for precise matching. The generated rule candidates are then normalized, deduplicated, and validated on a held-out corpus before being incorporated into the symbolic rule base. For example, if the neural extractor repeatedly maps \texttt{fabric.Connection(...).run(...)} to \texttt{shell\_interpreter\_execution}, the spec generator can synthesize a wrapper rule that binds \texttt{Connection.run} to the corresponding execution sink under lightweight context constraints.

\subsection{Skill Dependency Graph Generation}
\label{subsec:sdg_generation}

The second stage converts isolated SSO capabilities into a structured \emph{Skill Dependency Graph} (SDG) by resolving the logical operands that connect multiple operations. This stage consists of three components. First, \toolname performs \emph{symbolic operand resolution} to recover the security-relevant operands of each extracted SSO and trace their explicit origins through static point-to analysis. Second, it applies \emph{LLM-assisted operand inference} to normalize semantically equivalent operands and infer missing operand semantics that cannot be fully recovered through symbolic analysis alone. Third, it performs \emph{operation linking} by connecting SSOs that act on the same operand, or on operands related through derivation or control dependence, thereby constructing an operand-centric SDG. 

\begin{table}[t]
\centering
\small
\caption{Representative Operand Schemas.}
\label{tab:operand-types}
\begin{tabular}{ll}
\toprule
\textbf{Type} & \textbf{Operand Schema} \\
\midrule
\textbf{Exec} 
& \texttt{\{cmd, interpreter, argv, target\}} \\

\textbf{Net} 
& \texttt{\{endpoint, method, payload, response\}} \\

\textbf{File} 
& \texttt{\{path, mode, content\}} \\

\textbf{Env} 
& \texttt{\{name, value\}} \\

\textbf{Install} 
& \texttt{\{package, source, install\_cmd\}} \\

\textbf{Crypto} 
& \texttt{\{algorithm, key, input, output\}} \\

\bottomrule
\end{tabular}
\end{table}

\noindent \textbf{Symbolic Operand Resolution.} For each SSO record, we first recover the operand slots relevant to security reasoning according to the operation type. As summarized in \autoref{tab:operand-types}, each major SSO category is associated with a canonical operand schema. For example, execution operations expose command-oriented operands, file operations expose path- and content-oriented operands, and network operations expose endpoint- and payload-oriented operands. We then leverage an existing points-to analysis engine~\cite{yasa-engine,wang2026yasa} to determine the source of these operands and whether multiple operations refer to the same underlying object. Specifically, \toolname traces assignments, aliases, parameter passing, return values, wrapper calls, and inter-procedural data dependencies within and across files. This allows \toolname to recover explicit operand propagation chains, such as a secret read from \texttt{os.getenv(...)} being stored in a variable and later passed into \texttt{requests.post(...)}. 

\noindent \textbf{LLM-assisted Operand Inference.} Symbolic resolution is effective when operand propagation is explicit, but real-world skills are often incomplete, weakly typed, multi-file, or mixed with prompts, manifests, and setup instructions. As a result, some operand identities or semantic roles cannot be fully recovered through static analysis alone. To address this, we introduce an LLM-assisted operand inference component that normalizes semantically equivalent operands, infers missing operand semantics from local context and demonstrations, and bridges the gaps left by symbolic analysis. For instance, the model may infer that a variable named \texttt{token}, a prompt phrase such as ``API key,'' and a manifest field for credentials all denote the same secret-bearing operand class, even when no direct syntactic flow is visible. Likewise, it can normalize different textual descriptions of remote resources, shell commands, or installation targets into comparable operand representations. 

\noindent \textbf{SDG Construction.}  
After extracting SSO records and resolving their security-relevant operands, \toolname organizes them into a Security Dependency Graph (SDG), which captures how dispersed security-sensitive operations are related through structural containment, operand association, recovered value-flow dependencies, and shared operand references. Formally, we model the SDG as a typed directed graph \(G = (V, E, \phi_V, \phi_E)\), where \(\phi_V\) and \(\phi_E\) assign semantic types to nodes and edges, respectively. ~\autoref{fig:sdg-rules} summarizes the node and edge construction rules used to instantiate the SDG.

\begin{figure}[t]
\centering
\fbox{
\begin{minipage}{0.92\linewidth}
\footnotesize
\textbf{Notations:}
\begin{itemize}[leftmargin=15pt]
    \item \(a \in \mathcal{A}\): A program artifact, such as SKILL.md or other analysis unit.
    \item \(r \in \mathcal{R}\): A normalized SSO record.
    \item \(o \in \mathcal{O}\): A resolved operand object.
    \item \(\mathrm{artifact}(x)\): The artifact from which node \(x\) originates.
    \item \(\mathrm{opnds}(r)\): The set of operands associated with SSO record \(r\).
    \item \(\mathrm{resolve}(x)\): The operand object resolved from operand reference \(x\).
    \item \(o_1 \leadsto o_2\): A recovered value-flow dependency indicating that the value represented by \(o_1\) may propagate to \(o_2\).
\end{itemize}

\textbf{Node Types:}
\begin{enumerate}[leftmargin=15pt, label=\arabic*.]
    \item \textbf{Artifact Node:}
    \[
    a \in \mathcal{A}
    \quad \Rightarrow \quad
    \mathrm{type}(a) = \textsc{Artifact}
    \]

    \item \textbf{SSO Node:}
    \[
    r \in \mathcal{R}
    \quad \Rightarrow \quad
    \mathrm{type}(r) = \textsc{SSO}
    \]

    \item \textbf{Operand Node:}
    \[
    o \in \mathcal{O}
    \quad \Rightarrow \quad
    \mathrm{type}(o) = \textsc{Operand}
    \]
\end{enumerate}

\textbf{Edge Construction Rules:}
\begin{enumerate}[leftmargin=15pt, label=\arabic*.]
    \item \textbf{Containment Edge:}
    \[
    \mathrm{artifact}(r)=a
    \quad \Rightarrow \quad
    (a, \textit{contains}, r) \in E
    \]
    \[
    \mathrm{artifact}(o)=a
    \quad \Rightarrow \quad
    (a, \textit{contains}, o) \in E
    \]

    \item \textbf{Operand Association Edge:}
    \[
    x \in \mathrm{opnds}(r), \; \mathrm{resolve}(x)=o
    \quad \Rightarrow \quad
    (r, \textit{has\_operand}, o) \in E
    \]

    \item \textbf{Value-flow Dependency Edge:}
    \[
    o_1 \leadsto o_2
    \quad \Rightarrow \quad
    (o_1, \textit{flows\_to}, o_2) \in E
    \]

    \item \textbf{Shared-operand Edge:}
    \[
    \begin{aligned}
    &\exists o \in \mathcal{O}, \;
    (r_1, \textit{has\_operand}, o) \in E, \\
    &(r_2, \textit{has\_operand}, o) \in E, \;
    r_1 \neq r_2 \\
    &\Rightarrow
    (r_1, \textit{same\_operand}, r_2) \in E
    \end{aligned}
    \]
\end{enumerate}
\end{minipage}
}
\caption{Node and Edge Types in SDG Construction.}
\label{fig:sdg-rules}
\end{figure}

The node set \(V\) contains three categories. \emph{Artifact} nodes represent source-level program artifacts, such as files, modules, or other analysis units. \emph{SSO} nodes represent normalized security-sensitive operations extracted by \toolname. \emph{Operand} nodes represent resolved abstract operand objects corresponding to security-relevant operands referenced by SSO records. The edge set \(E\) contains four relation types. A \emph{contains} edge links an artifact node to an SSO node or operand node originating from that artifact. A \emph{has\_operand} edge links an SSO node to an operand node when that operand is explicitly referenced by the operation. A \emph{flows\_to} edge links two operand nodes when the value represented by one operand may propagate to the other through assignments, aliases, parameter passing, return-value binding, wrapper calls, or inter-procedural propagation. A \emph{same\_operand} edge links two SSO nodes when they reference the same resolved operand object after normalization and alias reconciliation; this is a derived correlation relation that makes operation-level sharing explicit in the graph.

The resulting SDG lifts isolated local observations into higher-level security-relevant behavior chains. For example, if a value obtained from \texttt{os.getenv(...)} is propagated through assignments and parameter passing into \texttt{requests.post(...)}, \toolname creates SSO nodes for the corresponding secret-access and network-transmission operations, operand nodes for the associated abstract values, and \emph{flows\_to} edges capturing the recovered propagation between operands. If both operations reference the same resolved operand object, the corresponding SSO nodes are additionally linked by a \emph{same\_operand} edge. This representation enables \toolname to reason about multi-step behaviors such as secret collection followed by exfiltration, suspicious command construction followed by execution, and file discovery followed by content access.


\subsection{Neuro-Symbolic Reasoning}
\label{subsec:neuro_symbolic_reasoning}

The third stage performs malicious-skill detection by combining formal Datalog reasoning with language-model-based reasoning in a neuro-symbolic loop. It takes the SDG produced in the previous stage and determines whether the observed operations and dependency structure instantiate known malicious patterns or suggest previously unseen ones.

\subsubsection{Fact Base Construction}
\label{subsubsec:fact_base}

We first convert the SDG into a relational \emph{fact base}. The fact base contains three categories of facts:
\begin{itemize}
    \item \textbf{Primitive operation facts}, which encode resolved security-sensitive actions such as environment reads, file writes, network sends, execution, dynamic loading, or package installation.
    \item \textbf{Operand relation facts}, which describe identity, derivation, and dependency relations among operands recovered during SDG generation.
    \item \textbf{Semantic annotation facts}, which specify security-relevant classes such as sensitive keys, external destinations, executable paths, code-like objects, startup-related locations, or externally controlled sources.
\end{itemize}

These facts form the extensional knowledge on which formal reasoning operates. Because all primitives are linked to explicit operands, the fact base preserves the object-centric semantics required to distinguish harmless capability co-occurrence from malicious data or control flows.

\subsubsection{Rule Base Encoding}
\label{subsubsec:rule_base}

On top of the fact base, we define a curated \emph{rule base} in Datalog. The rule base is designed offline and encodes reusable malicious behavior patterns as declarative logic rules. Rather than merely checking whether risky primitives co-occur in the same skill, these rules explicitly require that the relevant operations share the same operand or manipulate operands linked through derivation or dependency relations.

The rules are organized into multiple layers. The first layer consists of object-linking rules that compute transitive and semantic relations among operands, such as aliasing and derivation closure. The second layer derives intermediate semantic concepts from primitive facts, such as whether an object is secret-bearing, externally controlled, code-like, persistence-related, or destined for an external endpoint. The third layer encodes malicious behavior patterns such as secret exfiltration, downloader-installer behavior, remote code execution, dynamic sink injection, unauthorized persistence, and capability mismatch. The final layer maps one or more matched malicious patterns to final maliciousness labels and detection categories.

Because the rule base is explicit and declarative, each detection result remains inspectable and auditable: one can trace which primitive facts and which rules jointly contributed to the final conclusion.

\subsubsection{Datalog-Based Reasoning and Neural Feedback}
\label{subsubsec:datalog_and_feedback}

Given a fact base and a rule base, a Datalog engine performs least-fixed-point evaluation to derive all entailed conclusions. The outputs include triggered malicious-pattern labels, final maliciousness decisions, and derivation traces that function as auditable proof chains. This formal reasoning process provides the symbolic backbone of our system.

However, no finite rule base can enumerate all malicious behaviors encountered in practice. To address this limitation, we incorporate a complementary language-model-based reasoning component for \emph{unknown-pattern skills}, namely skills whose SDGs exhibit suspicious structure but do not confidently match existing symbolic rules.

This neural component serves two purposes. First, it provides online semantic assessment when symbolic reasoning is inconclusive but the skill exhibits non-trivial suspicious behavior. Second, when the model repeatedly identifies a high-value pattern across multiple samples, the pattern can be abstracted into reusable \emph{demo rules and patterns}, reviewed by analysts, and promoted into either the Datalog rule base or the capability-extraction heuristics. This yields a second neuro-to-symbolic feedback loop at the reasoning stage: neural analysis helps discover new malicious behaviors, while stable high-confidence findings are eventually codified into explicit symbolic rules.

The final output of the framework is a structured detection report containing the maliciousness decision, matched malicious-pattern categories, supporting facts, triggered rules, and the proof chain connecting low-level evidence to the final judgment.
